\section*{Chapitre \ref{ch:b3:asymmetric-synthesis} --- Synthèse asymétrique}

\begin{solution}{exo:b3:asymmetric-synthesis:1}
ee de 80\,\%~: $\mathrm{er} = 1.8/0.2 = 9$ (90\,:\,10). er de 99\,:\,1~: $\mathrm{ee} = 98/100 = 98$\,\%. ee de 50\,\%~: 75\,\% et 25\,\%.
\end{solution}

\begin{solution}{exo:b3:asymmetric-synthesis:2}
$(1840 - 60)/(1840 + 60) = 0.937$~: ee de 93.7\,\%, soit un er de $1840/60 = 31$.
\end{solution}

\begin{solution}{exo:b3:asymmetric-synthesis:3}
Priorités O $>$ éthyle $>$ H. Avec O en haut, l'éthyle en bas à gauche et H en bas à droite, la
face avant est \textit{Si}, la face arrière \textit{Re}. Un groupe méthyle ajouté par la \omterm{def:b3:asymmetric-synthesis:faces}{face Re}, avec
les priorités OH $>$ éthyle $>$ méthyle $>$ H dans le produit, donne le (\textit{R})-butan-2-ol.
\end{solution}

\begin{solution}{exo:b3:asymmetric-synthesis:4}
Éthanol~: \omterm{def:b3:asymmetric-synthesis:topicity}{énantiotopes} (échangés par le plan miroir de la molécule~; même
déplacement dans un solvant achiral). Butan-2-ol~: \omterm{def:b3:asymmetric-synthesis:topicity}{diastéréotopes} (voisins
d'un stéréocentre~; déplacements différents).
\end{solution}

\begin{solution}{exo:b3:asymmetric-synthesis:5}
er 199~: $RT\ln 199 = \qty{13.1}{kJ/mol}$ à \qty{298}{K}, \qty{8.6}{kJ/mol} à \qty{195}{K}. Un ee de 80\,\% à \qty{25}{\celsius} correspond à un er de 9,
$\Delta\Delta G^\ddagger = RT\ln 9 = \qty{5.45}{kJ/mol}$~; à \qty{-78}{\celsius}, $\mathrm{er} = \eu^{5450/(8.314 \times 195.15)} = 28.7$, soit un ee de 93\,\%.
\end{solution}

\begin{solution}{exo:b3:asymmetric-synthesis:6}
\omterm{def:b3:asymmetric-synthesis:ee}{Pureté optique} $12.0/15.0 = 80$\,\%~: 90\,\% de (\textit{S}) et 10\,\% de (\textit{R}). Un faible pouvoir rotatoire, les impuretés,
les effets de concentration et de solvant, et un comportement non linéaire
entre pouvoir rotatoire et composition rendent la mesure peu fiable~; la
chromatographie sépare et compte directement les énantiomères.
\end{solution}

\begin{solution}{exo:b3:asymmetric-synthesis:7}
$s = \ln(0.40 \times 0.10)/\ln(0.40 \times 1.90) = -3.22/-0.274 = 12$.
\end{solution}

\begin{solution}{exo:b3:asymmetric-synthesis:8}
Avec le L-($+$)-tartrate de diéthyle, le (2\textit{S},3\textit{S})-2,3-époxyhexan-1-ol~; avec le D-($-$)-tartrate de diéthyle,
l'énantiomère, (2\textit{R},3\textit{R}).
\end{solution}

\begin{solution}{exo:b3:asymmetric-synthesis:9}
L = phényle, M = méthyle, S = H. Pour l'aldéhyde (\textit{S}), le conformère réactif a le phényle
perpendiculaire au plan du carbonyle et l'hydrogène à côté de la
trajectoire du nucléophile, qui attaque en anti par rapport au phényle. La
détermination des configurations donne le
(2\textit{S},3\textit{S})-3-phénylbutan-2-ol comme diastéréoisomère majoritaire.
\end{solution}

\begin{solution}{exo:b3:asymmetric-synthesis:10}
Au départ, les deux énantiomères sont présents en quantités égales, si bien que les produits
se forment dans le rapport $k_{\mathrm{fast}} : k_{\mathrm{slow}} = s$~: $\mathrm{ee}_p = (s - 1)/(s + 1)$, 0.905 pour $s = 20$. À mesure que la réaction avance,
l'énantiomère rapide s'épuise et l'ee du produit diminue, tandis que l'ee
du substrat croît vers 1~: on peut toujours purifier le substrat en
allant plus loin~; pas le produit.
\end{solution}

\begin{solution}{exo:b3:asymmetric-synthesis:11}
Avec une racémisation rapide~: jusqu'à 100\,\% de rendement, $\mathrm{er} = 99$, ee de 98\,\%. Sans elle, à 50\,\% de
conversion~: rendement d'au plus 50\,\%, et ee du produit de 93\,\% (l'ee
du substrat est le même, 93\,\%).
\end{solution}

\begin{solution}{exo:b3:asymmetric-synthesis:12}
Lorsque les deux complexes s'interconvertissent plus vite qu'ils ne
réagissent, le rapport des produits vaut
$k_B[\mathrm B]/(k_A[\mathrm A]) = 600/10 = 60$~: le complexe minoritaire B donne 98\,\% du produit (ee de 97\,\%). La
sélectivité est fixée par la différence entre les deux états de
transition, et non par les \omterm{def:b3:partition-functions:boltzmann}{populations} des intermédiaires.
\end{solution}

\begin{solution}{pb:b3:asymmetric-synthesis:1}
\textbf{1.} Le carbone de l'alcène qui porte l'azote a trois substituants
différents et deux faces différentes~; l'hydrogène qui s'y ajoute fait du carbone $\alpha$ de l'acide
aminé un stéréocentre.

\textbf{2.} Le racémique~: les deux faces sont attaquées à la même vitesse.

\textbf{3.} Son oxygène de carbonyle se lie au rhodium en même temps que la
liaison C=C~: le substrat est maintenu sous forme de chélate, dans une
géométrie fixe que le ligand peut discriminer.

\textbf{4.} Elle construit une poche chirale autour du métal (de symétrie $C_2$)~: les deux façons de
fixer le substrat, par une face ou par l'autre, sont diastéréoisomères,
d'énergies et de réactivités différentes.

\textbf{5.} Non~: le complexe minoritaire réagit bien plus vite avec le dihydrogène (Curtin--Hammett).

\textbf{6.} Le métal se trouve sur une face de l'alcène lié, et les
hydrogènes sont transférés depuis le métal.

\textbf{7.} $\mathrm{er} = 1.95/0.05 = 39$.

\textbf{8.} $\Delta\Delta G^\ddagger = RT\ln 39 = 8.314 \times 298.15 \times 3.664 = \qty{9.1}{kJ/mol}$.

\textbf{9.} $\mathrm{er} = \eu^{9082/(8.314 \times 195.15)} = 270$~: ee de 99.3\,\%.

\textbf{10.} $\mathrm{er} = \eu^{11082/(8.314 \times 298.15)} = 87$~: ee de 97.7\,\%.

\textbf{11.} Moins qu'une liaison hydrogène typique~: une petite différence
de contacts stériques ou une seule interaction faible entre le substrat et
le ligand décide du résultat, ce qui explique que de tels catalyseurs
soient difficiles à concevoir à partir des premiers principes.

\textbf{12.} La vitesse diminue, la solubilité du dihydrogène et le
comportement du catalyseur changent à basse température, et refroidir un
grand réacteur coûte de l'énergie.

\textbf{13.} 97.5 g de l'énantiomère (\textit{S}) et 2.5 g de l'énantiomère (\textit{R}).

\textbf{14.} $97.5 - 5.0 = \qty{92.5}{g}$ de cristaux, pratiquement énantiopurs.

\textbf{15.} $92.5/100 = 92.5$\,\% du produit brut (95\,\% de l'énantiomère (\textit{S})).

\textbf{16.} Les cristaux recueillent l'énantiomère en excès, la partie
racémique restant en solution avec l'énantiomère minoritaire~; à partir
d'un racémique, il n'y a aucun excès à recueillir, et la cristallisation
donne des cristaux racémiques.

\textbf{17.} Par HPLC chirale, avec une référence racémique.

\textbf{18.} 50\,\%~: seul l'énantiomère lent subsiste.

\textbf{19.} $s = \ln(0.45 \times 0.05)/\ln(0.45 \times 1.95) = -3.79/-0.131 = 29$.

\textbf{20.} 45\,\% du racémique (88\,\% de l'énantiomère voulu présent au
départ, avec un ee de 95\,\%).

\textbf{21.} Avec une racémisation rapide du substrat, le rendement pourrait approcher 100\,\%.

\textbf{22.} Elle convertit tout le précurseur achiral, bon marché, en le
bon énantiomère, avec un catalyseur utilisé en quantités infimes et sans
jeter d'énantiomère.

\textbf{23.} Le rhodium est toxique et ses résidus dans les médicaments sont
limités à quelques parties par million~; on les dose par spectrométrie de
masse à plasma à couplage inductif.

\textbf{24.} Un ee de 95\,\% à \qty{25}{\celsius} correspond à $\Delta\Delta G^\ddagger = RT\ln 39 \approx \qty{9.1}{kJ/mol}$.
\end{solution}
